\documentclass[10pt,a4paper]{amsart}
\usepackage[margin=19mm]{geometry}
\usepackage{amsmath,amssymb,mathtools}
\usepackage[scr=rsfs,cal=euler]{mathalfa}
\usepackage{xfrac}
\usepackage[unicode,hidelinks,bookmarks=false]{hyperref}
\newcommand{\ie}{i.e.\ }
\newcommand{\cf}{cf.\ }
\newcommand{\resp}{respectively\ }
\newcommand{\Subsection}{\subsection}
\providecommand{\nobreakdash}{\nobreak\hbox{-}}
\setlength{\emergencystretch}{2em}
\multlinegap0pt
\usepackage{amssymb}
\usepackage{mathtools}
\newtheorem{lemma}{Lemma}
\newcommand{\Exs}{\ensuremath{{\mathbb{E}}}}
\newcommand{\Prob}{\ensuremath{{\mathbb{P}}}}
\newcommand{\To}{{T_o}}
\newcommand{\abss}[1]{\left| #1 \right |}
\newcommand{\bdelta}{\bar{\delta}}
\newcommand{\braces}[1]{\left\{ #1 \right \}}
\newcommand{\brackets}[1]{\left[ #1 \right]}
\newcommand{\btau}{\kappa}
\newcommand{\cG}{\mathcal{G}}
\newcommand{\cM}{\mathcal{M}}
\newcommand{\cS}{\mathcal{S}}
\newcommand{\defn}{:=}
\newcommand{\excess}{\alpha}
\newcommand{\flip}{\varsigma}
\newcommand{\jL}{\bar{L}}
\newcommand{\parenth}[1]{\left( #1 \right)}
\newcommand{\rV}{\mathscr{V}}
\newcommand{\real}{\ensuremath{\mathbb{R}}}
\newcommand{\sF}{\mathscr{F}}
\newcommand{\stopT}{{\boldsymbol{\mathfrak{T}}}}
\newcommand{\tagi}{{[i]}}
\newcommand{\uV}{V}
\newcommand{\uW}{W}
\title{Talagrand's convolution conjecture up to $\log \log$ \\ via perturbed reverse heat}
\author{Yuansi Chen}
\date{Source: arXiv:2511.19374v2 (CC BY 4.0)}
\begin{document}
\maketitle

\noindent\emph{Source and licence.} Talagrand's convolution conjecture up to $\log \log$ \\ via perturbed reverse heat. Version arXiv:2511.19374v2 (CC BY 4.0), distributed by arXiv under the Creative Commons Attribution 4.0 International licence, http://creativecommons.org/licenses/by/4.0/, as recorded in the arXiv licence metadata for this exact version and shown on the abstract page https://arxiv.org/abs/2511.19374v2. Full licence notice: \texttt{licenses/arxiv-2511.19374-CC-BY-4.0.txt}.\par\smallskip

\noindent\textit{Editorial scope.} Counted unit: the article's lemma lem:weighted\_energy\_estimate (source0.tex lines 820--832). The article's complete original proof of it is delivered with it (source0.tex lines 1432--1499, one \texttt{proof} environment of 68 lines, which the article places away from the statement). No mathematics is written, repaired or re-derived: the statement and the proof are the article's own lines, in the article's own notation, and no retained line is substituted or re-flowed.  Retained uncounted as the source-local context the proof needs: lines 461--465; lines 479--482; lines 731--737; lines 741--744; lines 745--755; lines 758--785; lines 788--811.  Nothing was omitted from any retained span, so the package carries no omission record.  No retained line contains a citation, so the wrapper carries no bibliography and no bibliography entry.  No retained line contains a figure environment, an image-inclusion command or drawing code, so no image is delivered and the package declares no asset dependency.  Editorial formatting only, in addition: an AMS \texttt{amsart} preamble that restates the article's own declarations and macros used by the retained lines, namely the environments \texttt{lemma} on separate counters, together with the standard packages those lines need.  The retained environments print in this document as Lemma 1 in order of appearance, whereas the article prints the same items with the numbering of its own full document.  The complete native source of the article as distributed in the arXiv e-print for this exact version is bundled unchanged as \texttt{sources/189716/source0.tex}.
\begin{align}
  \label{eq:def_reverse_process}
  \uV_0 &\sim \nu_{P_T f}, \notag \\
  d \uV_t &=  \sum_{i=1}^n (-2 \uV_{t-} e_i) \int_\real \mathbf{1}_{0 < z \leq \frac12 - S_i(e^{-{(T-t)}} \uV_{t-})} N^\tagi(dt, dz),
\end{align}

\begin{align}
  \label{eq:def_scaled_reverse_jump_process}
  \rV_t \defn e^{-(T-t)} \uV_t, 
\end{align}

\begin{lemma}[Basic level-1 inequality]
  \label{lem:level_1_ineq}
  Let $h: \braces{-1,1}^n \to \braces{0, 1}$. Consider its multilinear expansion. For any $x \in (-1,1)^n$, we have
  \begin{align*}
    \sum_{i=1}^n (1-x_i^2) \parenth{\partial_i h(x)}^2 \leq h(x) - h(x)^2 \leq \frac14. 
  \end{align*}
\end{lemma}

\begin{align}
  \label{eq:def_cS}
  \cS \defn \Exs \brackets{ \int_{\theta}^{\stopT} \sum_{i=1}^n \bdelta^2 S_i(\rV_{t-})^2 dt}
\end{align}

\begin{lemma}[Expected squared score bound] 
  \label{lem:expected_squared_score_bound}
  For the reverse heat process~\eqref{eq:def_reverse_process} and its scaled version $(\rV_t)$ in Eq.~\eqref{eq:def_scaled_reverse_jump_process}, we have
  \begin{align*}
    \Exs \brackets{ \int_{\theta}^{\stopT} \sum_{i=1}^n S_i(\rV_{t-})^2 dt \middle| \sF_\theta} \leq \btau (R_\theta + \excess + 1).
  \end{align*}
  Consequently, $\cS$ defined in~\eqref{eq:def_cS} satisfies
  \begin{align*}
    \cS \leq \Exs \brackets{\parenth{\frac{\excess \mathbf{1}_{R_\theta \geq \excess}}{R_\theta + 1} }^2 \btau (R_\theta + \excess + 1)} \leq 2 \btau \excess^2 \Exs\brackets{\frac{\mathbf{1}_{R_\theta \geq \excess}}{R_\theta + 1}}.
  \end{align*}
\end{lemma}

\begin{lemma}[Boolean heat bridge]
  \label{lem:boolean_bridge_formula}
  For $0 \leq t \leq \To$ and $x, y, \zeta \in \braces{-1,1}^n$, define
\begin{align}
  \label{eq:def_m}
  m_{t}(x, y, \zeta) \defn \Exs\brackets{\uV_{\To} \odot (x \odot y) \middle| \uV_t = x, \uV_{T} = \zeta}.
\end{align}
The coordinates of $\uV_{\To}$ are conditionally independent given $(\uV_t, \uV_{T})$ and their conditional law has a closed form. For $i \in [n]$, the $i$-th component of $m_{t}(x, y, \zeta)$ is given by  
  \begin{align*}
     m_{t}^{[i]}(x, y, \zeta) = a_{t} y_i + b_{t} \sigma_i
  \end{align*}
  where $\sigma_i \defn x_i y_i \zeta_i$ and
  \begin{align}
    \label{eq:def_a_b}
    a_{t} &\defn \frac{\sinh(T-\To)}{\sinh(T-t)}, \quad b_{t} \defn \frac{\sinh(\To-t)}{\sinh(T-t)}.
  \end{align}
  Because the coordinates of $\uV_{\To}$ are conditionally independent, for any multilinear $\phi:\braces{-1,1}^n \to \real$,
  \begin{align*}
    \Exs\brackets{\phi(\uV_\To \odot (x \odot y))  \middle| \uV_t = x, \uV_T = \zeta} = \phi(m_{t}(x, y, \zeta)).
  \end{align*}
  We define $q_t^{\zeta}(x, y) \defn \phi(m_{t}(x, y, \zeta))$. Then
  \begin{align}
    \label{eq:q_t_derivative}
    \Delta_i^y q_t^{\zeta}(x, y) &= -2 (a_t y_i + b_t \sigma_i) \partial_i \phi(m_{t}(x, y, \zeta)), \notag \\
    \Delta_i^y q_t^{\zeta}(\flip_i(x), y) &= -2 (a_t y_i - b_t \sigma_i) \partial_i \phi(m_{t}(x, y, \zeta)), \notag \\
    \Delta_{i}^{xy} q_t^{\zeta}(x, y) &= -2 a_t y_i \partial_i \phi(m_{t}(x, y, \zeta)).
  \end{align}
\end{lemma}

\begin{lemma}[Doob $h$-transform of the predictable joint generator]
  \label{lem:Doob_h_transform}
  Given $\zeta \in \braces{-1,1}^n$, define
  \begin{align}
    \label{eq:def_H_r_lambda}
    H_t^\zeta(x) &\defn \Prob(\uV_T = \zeta \mid \uV_t = x), \notag \\
    r_{t, i}^{\zeta}(x) &\defn \frac{H_t^\zeta(\flip_i(x))}{H_t^\zeta(x)}, \notag \\
    \lambda_{t, i}^{\zeta}(x) &\defn \frac{1 - \rho_t x_i \zeta_i}{1 + \rho_t x_i \zeta_i}, \quad \forall x \in \braces{-1, 1}^n.
  \end{align}
  Under the conditioned law $\Prob^\zeta \defn \Prob(\cdot \mid \uV_T = \zeta)$,  with shorthand $\rho_t = e^{-(T-t)}$, $\delta_i = \delta_i(\rho_t x)$, $S_i = S_i(\rho_t x)$, $r_i^{\zeta} = r_{t, i}^{\zeta}(x)$, $\lambda_i^{\zeta} = \lambda_{t, i}^{\zeta}(x)$, for every bounded function $h: \braces{-1, 1}^{2n} \to \real$, 
  \begin{align*}
    \cM_t^{\zeta, h} &\defn h(\uV_t, \uW_t) - h(\uV_\theta, \uW_\theta) - \int_\theta^t \cG_s^\zeta h(\uV_{s-}, \uW_{s-}) ds
  \end{align*}
  is a $\Prob^\zeta$-martingale on $[\theta, T]$, where $\cG_t^\zeta$ is the Doob $h$-transform of the predictable joint generator $\cG_t$ with respect to the function $H_t^\zeta$, and it takes the form
  \begin{align*}
    \cG_t^\zeta &= \jL_t^{0, \zeta} + \mathbf{1}_{t \leq \stopT} A_t^\zeta, \\
    \jL_t^{0, \zeta} h(x, y) &= \sum_{i=1}^n r_{i}^{\zeta} \parenth{\frac12 - S_i } \Delta_i^{xy} h(x, y) \\
    A_t^\zeta h(x, y) &= \sum_{i=1}^n  \mathbf{1}_{S_i > 0}  \delta_i  S_i \Delta_{i}^y h(x, y)  + \mathbf{1}_{S_i \leq 0} r_{i}^{\zeta} \delta_i  S_i  \Delta_i^y h(\flip_i(x), y).
  \end{align*}
  Moreover,
  \begin{align*}
    r_i^{\zeta} \parenth{\frac12 - S_i } &=  \frac12 \lambda_{i}^\zeta.
  \end{align*}
\end{lemma}

\begin{lemma}[Weighted energy estimate]
  \label{lem:weighted_energy_estimate}
  Given a function $\phi:\braces{-1,1}^n \to \braces{0, 1}$. For $a_{t}$, $b_{t}$ and $m_{t}$ defined in Lemma~\ref{lem:boolean_bridge_formula}, $\lambda_{t, i}^{\zeta}(x)$ defined in Lemma~\ref{lem:Doob_h_transform}, define
  \begin{align*}
    \Psi_a &\defn \Exs \int_\theta^{\To} a_{t}^2 \sum_{i=1}^n \lambda_{t, i}^{\uV_T}(\uV_t) \abss{\partial_i \phi(m_{t}(\uV_t, \uW_t, \uV_T))}^2 dt, \\
    \Psi_b &\defn \Exs \int_\theta^{\To} b_{t}^2 \sum_{i=1}^n \lambda_{t, i}^{\uV_T}(\uV_t) \abss{\partial_i \phi(m_{t}(\uV_t, \uW_t, \uV_T))}^2 dt.
  \end{align*}
  Then, with $\cS$ defined in Lemma~\ref{lem:expected_squared_score_bound}, we have
  \begin{align*}
    \Psi_b &\leq \frac{e^{-2\tau}}{1-e^{-2\tau}} \parenth{\To - \theta}, \\
    \Psi_a &\lesssim 1 + \btau \cS + \Psi_b.
  \end{align*}
\end{lemma}

\begin{proof}[Proof of Lemma~\ref{lem:weighted_energy_estimate}]
  First, we bound $\Psi_b$. For $t \in [\theta, \To]$, $x, y, \zeta \in \braces{-1, 1}^n$ and $i \in [n]$, the definition of $m_t$ in Lemma~\ref{lem:boolean_bridge_formula} yields 
  \begin{align*}
    1- m_t^{[i]}(x, y, \zeta)^2 &= \frac{(1-\rho_{\To}^2)(\rho_{\To}^2 - \rho_t^2)} {\rho_{\To}^2 (1+\rho_t x_i \zeta_i)^2} \\
    b_t^2 &= \frac{\sinh^2(\To-t)}{\sinh^2(T-t)} = \frac{(\rho_{\To}^2 - \rho_t^2)^2}{\rho_{\To}^2 (1-\rho_t^2)^2}.
  \end{align*}
  Taking the ratio and multiplying by $\lambda_{t, i}^{\zeta}(x)$ yields, for $t < \To$,
  \begin{align*}
    \frac{\lambda_{t, i}^{\zeta}(x) b_t^2}{1- m_t^{[i]}(x, y, \zeta)^2} = \frac{ \rho_{\To}^2 - \rho_t^2}{1-\rho_t^2} \cdot \frac{1}{1-\rho_{\To}^2} \leq \frac{\rho_{\To}^2}{1-\rho_{\To}^2} = \frac{e^{-2\tau}}{1-e^{-2\tau}}.
  \end{align*}
  Hence,
  \begin{align*}
    \lambda_{t, i}^{\zeta}(x) b_t^2 \abss{\partial_i \phi(m_t(x, y, \zeta))}^2 &\leq \frac{e^{-2\tau}}{1-e^{-2\tau}} \parenth{1- m_t^{[i]}(x, y, \zeta)^2} \abss{\partial_i \phi(m_t(x, y, \zeta))}^2.
  \end{align*}
  Summing over $i$, applying the level-1 inequality in Lemma~\ref{lem:level_1_ineq} and integrating, we conclude that $\Psi_b \leq \frac{e^{-2\tau}}{1-e^{-2\tau}} (\To - \theta)$.

  Second, we bound $\Psi_a$. Fix $\zeta \in \braces{-1,1}^n$ and work under the conditioned law $\Prob^\zeta$. From $m_t$ and $q_t^{\zeta}$ in Lemma~\ref{lem:boolean_bridge_formula}, we have
  \begin{align*}
    \partial_t m_t^{[i]}(x, y, \zeta) &= \partial_t a_t y_i + \partial_t b_t \sigma_i = a_t y_i \frac{1-\rho_t x_i \zeta_i}{1+\rho_t x_i \zeta_i} \\
    \Delta_i^{xy} q_t^{\zeta}(x, y) & = - 2 a_t y_i \partial_i \phi(m_t(x, y, \zeta)).
  \end{align*}
  Using the multilinearity of $\phi$ and the unperturbed generator representation in Lemma~\ref{lem:Doob_h_transform}, we verify that
  \begin{align}
    \label{eq:q_zeta_harmonic}
    \parenth{\partial_t + \jL_t^{0, \zeta}} q_t^{\zeta} = 0.
  \end{align}
  Using the conditioned predictable joint generator from
  Lemma~\ref{lem:Doob_h_transform}, apply It\^o's formula to $q_{t}^{\zeta}(\uV_t, \uW_t)^2$ under the conditioned law $\Prob^\zeta$ to obtain
  \begin{align}
    \label{eq:q_zeta_sq_ito}
    \Exs^{\zeta} \brackets{q_{\To}^{\zeta}(\uV_{\To}, \uW_{\To})^2 - q_{\theta}^{\zeta}(\uV_\theta, \uW_\theta)^2} &= \Exs^{\zeta} \int_\theta^{\To} \parenth{\partial_t + \jL_t^{0, \zeta} + \mathbf{1}_{t \leq \stopT} A_t^{\zeta}} \parenth{q_{t}^{\zeta}}^2 (\uV_{t-}, \uW_{t-}) dt.
  \end{align}
  For the unperturbed part, Eq.~\eqref{eq:q_zeta_harmonic} yields
  \begin{align*}
    \parenth{\partial_t + \jL_t^{0, \zeta}} \parenth{q_{t}^{\zeta}}^2 (x, y) &= \sum_{i=1}^n \frac12 \lambda_{t, i}^{\zeta}(x) \parenth{\Delta_i^{xy} q_{t}^{\zeta} (x, y)}^2 \\
    &= 2 a_t^2\sum_{i} \lambda_{t, i}^{\zeta}(x) \abss{\partial_i \phi(m_t(x, y, \zeta))}^2.
  \end{align*}
  For the perturbation part, recall that 
  \begin{align*}
    A_t^\zeta h(x, y) &= \sum_{i=1}^n  \mathbf{1}_{S_i > 0}  \delta_i  S_i \Delta_{i}^y h(x, y)  + \mathbf{1}_{S_i \leq 0} r_{i}^{\zeta} \delta_i  S_i  \Delta_i^y h(\flip_i(x), y).
  \end{align*}
  Note that from Lemma~\ref{lem:Doob_h_transform} and $\btau^{-1} \leq \lambda_{t, i}^{\zeta}(x) \leq \btau$,
  \begin{align*}
    \abss{\delta_i \brackets{\mathbf{1}_{S_i > 0} + r_i^{\zeta} \mathbf{1}_{S_i \leq 0} }}^2 = \abss{\bdelta \brackets{\mathbf{1}_{S_i > 0} + \frac{\lambda_{t, i}^{\zeta}(x)}{1 - 2 \bdelta S_i} \mathbf{1}_{S_i \leq 0}}}^2
    \leq \bdelta^2 \btau  \lambda_{t, i}^{\zeta}(x).
  \end{align*}
  Additionally, since $q^{\zeta}$ is bounded by $1$ because $\phi$ as a multilinear extension takes value in $[0, 1]$, we have
  \begin{align*}
    \abss{\Delta_i^y (q_{t}^{\zeta})^2 (x, y)} &\leq 2 \abss{\Delta_i^y q_{t}^{\zeta} (x, y)} \\
    \abss{\Delta_i^y (q_{t}^{\zeta})^2 (\flip_i(x), y)} &\leq 2 \abss{\Delta_i^y q_{t}^{\zeta} (\flip_i(x), y)}.
  \end{align*}
  Hence,
  \begin{align*}
    \abss{A_t^{\zeta} (q_{t}^{\zeta})^2 (x, y)} &\lesssim \bdelta \btau^{\frac12} \sum_{i=1}^n \parenth{\lambda_{t, i}^{\zeta}(x)}^{\frac12} \abss{S_i} \parenth{\abss{\Delta_i^y q_{t}^{\zeta} (x, y)} + \abss{\Delta_i^y q_{t}^{\zeta} (\flip_i(x), y)}} \\
    &\overset{(i)}{\lesssim} \btau^{\frac12} \parenth{\sum_{i=1}^n \bdelta^2 S_i^2}^{\frac12} \parenth{\sum_{i=1}^n \lambda_{t, i}^{\zeta}(x) \parenth{\abss{\Delta_i^y q_{t}^{\zeta} (x, y)}^2 + \abss{\Delta_i^y q_{t}^{\zeta} (\flip_i(x), y)}^2}}^{\frac12} \\
    &\overset{(ii)}{\lesssim} \btau^{\frac12} \parenth{\sum_{i=1}^n \bdelta^2 S_i^2}^{\frac12} \parenth{\sum_{i=1}^n \lambda_{t, i}^{\zeta}(x) \parenth{a_t^2 + b_t^2}\abss{\partial_i \phi(m_t(x, y, \zeta))}^2 }^{\frac12}
  \end{align*}
  where (i) follows from Cauchy-Schwarz inequality, and (ii) follows from $\Delta_i^y q_{t}^{\zeta}$ expression in Lemma~\ref{lem:boolean_bridge_formula}. 
  Evaluating at $(\uV_{t-}, \uW_{t-})$, integrating in time, applying Cauchy-Schwarz, averaging over $\zeta = \uV_T$ and plugging back in Eq.~\eqref{eq:q_zeta_sq_ito} yields
  \begin{align*}
    \Psi_a &\lesssim 1 + \btau^{\frac12} \cS^{\frac12} \parenth{\Psi_a + \Psi_b}^{\frac12} \\
    &\lesssim 1 + \btau \cS + \frac14 (\Psi_a + \Psi_b).
  \end{align*}
  Absorbing $\frac14 \Psi_a$ to the left-hand side yields the desired bound
  \begin{align*}
    \Psi_a &\lesssim 1 + \btau \cS + \Psi_b.
  \end{align*}
\end{proof}

\end{document}
